% !PRINTED 431
\weberpage{431}
%<<CONT
\gesperrt{befriedigt wird. Dieselbe Congruenz wird aber auch befriedigt, wenn als
Exponent eine mit $\alpha$ nach dem Modul $n-1$ congruente Zahl gesetzt
wird, und man findet in jedem vollen Restsystem nur eine solche Zahl
$\alpha$.}

Der Exponent $\alpha$ heisst der \gesperrt{Index} von $a$ in Bezug auf die \gesperrt{Basis}
$g$, und es wird geschrieben
\weq{17}{\alpha \equiv \operatorname{ind} a \pmod{n-1}.}

Man hat also für jede durch $n$ nicht theilbare Zahl
\weq{18}{g^{\operatorname{ind} a} \equiv a \pmod{n}.}

Aus den beiden Formeln
\weq{19}{g^{\operatorname{ind} a + \operatorname{ind} b} \equiv ab, \quad g^{m \operatorname{ind} a} \equiv a^m \pmod{n}}
ergiebt sich der Satz:

6.\ \gesperrt{Der Index eines Productes ist gleich der Summe der Indices
der Factoren; der Index der $m^{\mathrm{ten}}$ Potenz von $a$ ist gleich
dem $m$fachen des Index von $a$.}

Diese Sätze sind ganz analog den entsprechenden Sätzen, die sich auf die
Rechnung mit Logarithmen beziehen, und die Analogie lässt sich noch
weiter verfolgen. So wird z.\,B. der Uebergang von einer Basis $g$ zu
einer anderen Basis $g'$ durch die Formel
\weq{20}{\overset{g}{\operatorname{ind}}\, a \equiv \overset{g}{\operatorname{ind}}\, g' \; \overset{g'}{\operatorname{ind}}\, a \pmod{n-1}}
vermittelt. Wir wollen noch den leicht zu beweisenden Satz anführen,
dass unter den durch $n$ nicht theilbaren Zahlen $a$ \gesperrt{die und nur die}
primitive Wurzeln sind, deren Indices relativ prim zu $n-1$ sind.

Für praktische Rechnungen bedient man sich zweckmässig sogenannter
Indextabellen, die den Logarithmentafeln entsprechen.

Man nennt in der Congruenz (17) $\alpha$ den Index und $a$ den Numerus,
und stellt am besten zwei Tabellen auf, von denen die eine zu jedem
Index den Numerus, die andere zu jedem Numerus den Index giebt, wo die
zweite durch Umstellung aus der ersten gewonnen wird. Dabei ist zu
beachten, dass für die Indices der Modul $n-1$, für die Numeri der
Modul $n$ in Betracht kommt und congruente Zahlen als gleichwerthig
gelten. So erhält man z.\,B. für $n = 13$ und die Basis 2
%CONT>>
